sending @@M0@@ to @@M1@@, and @@M2@@ to @@M3@@. It follows immediately from (5.5.1) and (5.3.1) that this is a @@M4@@-homomorphism, and, since @@M5@@ is an integral domain in the case considered, one immediately verifies that it is injective; since @@M6@@ is a @@M7@@-ring, we deduce the result.

In the @@M8@@-ring @@M9@@, we can define operations @@M10@@ (3.2). For @@M11@@, the value of @@M12@@ calculated on the element @@M13@@ of @@M14@@ is given by:
@@M15@@
it is therefore an element of @@M16@@; it is defined by the relation:
\sgaStableTarget{sga6.e05.eq.5-5-2}@@M17@@
and the universal polynomials resulting from it.

\sgaStableTarget{sga6.e05.para.5-6}5.6. Let @@M18@@ be the @@M19@@-ring generated over @@M20@@ by the generators @@M21@@, subject to the relations @@M22@@ for @@M23@@, @@M24@@ for @@M25@@. We propose to calculate the @@M26@@ in terms of the @@M27@@.

By definition:
@@M28@@
that is:
@@M29@@

